Tissue function depends on spatial organization across scales, from cellular neighborhoods to anatomical regions. Expanding spatial context around a focal location yields related but distinct representations of tissue organization at different granularities. Although existing spatial transcriptomics pretraining models incorporate multiscale context for representation learning, they do not explicitly learn how representations change as spatial context expands. We introduce ZoomST, a self-supervised framework that jointly learns granularity-specific, cross-scale-compatible representations and transformations among them. ZoomST combines molecular-spatial dual-view alignment within each granularity with a soft distributional equivariance constraint across granularities. Across held-out mouse sections spanning all sampled developmental stages, frozen ZoomST representations outperform the evaluated pretrained baselines and remain competitive with target-fitted methods for tissue mapping. Anatomical analysis in a mouse developmental atlas further illustrates the complementary information captured by fine and coarse spatial context. Continuous transformations learned only from discrete training granularities predict representations at unseen intermediate spatial scales and reveal relative differences in context sensitivity across tissue locations. By making spatial granularity adjustable after pretraining, ZoomST turns changes in spatial context into both a source of self-supervision and an axis for tissue analysis.
